\documentclass[11pt,oneside]{article}
\pdfminorversion=7
\usepackage{CourseNotesStyle}
\hypersetup{
  pdftitle={Course Notes for CSED490V Foundations of Continuous Computation},
  pdfauthor={Jungwoo Kim},
  pdfsubject={Computable real numbers and computable functions},
  pdfcreator={LaTeX}
}

\begin{document}
\thispagestyle{empty}
\begin{center}
  \vspace*{1.1cm}
  {\LARGE\bfseries Course Notes\par}
  \vspace{0.25cm}
  {\large for (CSED490V) Foundations of Continuous Computation\par}
  \vspace{0.75cm}
  {\normalsize Jungwoo Kim\\
   Department of Mathematics\\
   \PohangUniversityOfScienceAndTechnology\par}
\end{center}
\vspace{0.45cm}

These notes develop the mathematical foundations of continuous computation. General
mathematical notation and writing follow \textit{(MATH000) Mathematical Language}.

\tableofcontents

\section*{Notation and Conventions}
\phantomsection\addcontentsline{toc}{section}{Notation and Conventions}
Only notation specific to this course is listed here. Its mathematical meaning is
made precise at the point of first use.
\begin{coursenotation}
$\Rc$ & Set of computable real numbers\\
$\A$ & An algorithm, possibly with access to a real-number name
\end{coursenotation}

\section{Computable Reals}

A real number is computable when it can be approximated by rational numbers to
any prescribed precision through a finite computation. There are two useful
ways to express this idea: rational approximation and binary expansion.

\begin{definition}[Borel Computability]
A real number $r$ is \emph{Borel-computable} if there exists a computable
function $q\colon\Nzero\to\Q$ such that
\[
  |r-q(n)|<2^{-n}\qquad(n\in\Nzero).
\]
Equivalently, an algorithm computes a rational approximation to $r$ within
$2^{-n}$ when given a precision parameter $n$.
\end{definition}

\begin{definition}[Turing Computability]
A nonnegative real number $r$ is \emph{Turing-computable} if it has a binary
expansion
\[
 r=a+\sum_{i=0}^{\infty}b_i2^{-(i+1)},
 \qquad a\in\Nzero,\quad b_i\in\{0,1\},
\]
for which $i\mapsto b_i$ is computable. A negative real number is
Turing-computable when its absolute value is Turing-computable. For dyadic
rationals we may use the terminating binary expansion.
\end{definition}

\begin{theorem}[Equivalence of Computability Definitions]
For $r\in\R$, the following are equivalent:
\begin{enumerate}[label=(\roman*)]
 \item $r$ is Borel-computable.
 \item $r$ is Turing-computable.
\end{enumerate}
\end{theorem}
\begin{proof}
We may consider $r\ge0$: changing the sign of a computable rational
approximation is effective, and the negative case of binary computability is
defined using the absolute value.

Suppose first that $r$ is Turing-computable, with integer part $a$ and
computable bits $b_i$. Define
\[
 q(n)=a+\sum_{i=0}^{n}b_i2^{-(i+1)}\in\Q.
\]
The sum is computable and
\[
 0\le r-q(n)\le\sum_{i=n+1}^{\infty}2^{-(i+1)}
   =2^{-(n+1)}<2^{-n}.
\]
Thus $r$ is Borel-computable.

Conversely, suppose $q\colon\Nzero\to\Q$ computes $r$ with
$|q(m)-r|<2^{-m}$. Fix the integer part $a=\lfloor r\rfloor$ as part of the
existential description of a binary expansion. We distinguish two cases.

If $r$ is dyadic, it admits an eventually zero binary expansion. Its finite
sequence of nonzero digits can be built into an algorithm, and hence $r$ is
Turing-computable.

Now assume $r$ is not dyadic. Fix $n\in\Nzero$. The non-dyadic number $r$
lies strictly between two consecutive grid points of mesh $2^{-(n+1)}$.
Search over $m\in\Nzero$ until the rational interval
\[
  I_m=(q(m)-2^{-m},\,q(m)+2^{-m})
\]
is contained strictly in a single grid cell
\[
 \left(\frac{k}{2^{n+1}},\frac{k+1}{2^{n+1}}\right),
 \qquad k\in\Nzero.
\]
Containment is decidable by rational arithmetic: for a given $m$ there is
at most one candidate $k$, which can be found by integer division.
The search terminates, since $r$ has positive distance from both endpoints
of its grid cell and $q(m)\to r$ with an effective error bound.
Moreover $r\in I_m$, so the returned $k$ satisfies
$k=\lfloor2^{n+1}r\rfloor$. The $n$th binary digit of $r$ is therefore
\[
 b_n=k\bmod 2.
\]
This procedure computes each bit, completing the proof.
\end{proof}

\begin{remark}
The two definitions characterize the same set of real numbers, but their
representations do not admit symmetric \emph{uniform} conversions. Given a
computable binary expansion one can effectively obtain rational approximations.
In general there is no algorithm that takes an arbitrary rational-approximation
procedure and uniformly produces its binary digits: the latter operation must
resolve boundary cases at dyadic rationals. This does not contradict the
existence of a computable binary expansion for each Borel-computable real.
\end{remark}

\begin{definition}[Computable Real Numbers]
The set of \emph{computable real numbers} is
\[
  \Rc:=\{r\in\R:r\text{ is computable}\}.
\]
By the preceding theorem, either notion of computability can be used.
\end{definition}

\begin{proposition}[Closure under Arithmetic]
The set $\Rc$ is a subfield of $\R$: if $x,y\in\Rc$, then
\[
 x+y,\quad x-y,\quad xy\in\Rc,
 \qquad \frac{x}{y}\in\Rc\quad\text{when }y\ne0.
\]
\end{proposition}
\begin{proof}
Take computable rational approximations to $x$ and $y$ with arbitrarily
small certified error. Addition and subtraction preserve effective error
bounds by the triangle inequality. For multiplication, first obtain rational
upper bounds on $|x|$ and $|y|$, then use
\[
 |xy-uv|\le |x|\,|y-v|+|v|\,|x-u|
\]
to choose sufficiently accurate rational $u,v$. Finally, if $y\ne0$,
search for an approximation $v$ whose certified error is smaller than $|v|$;
this eventually succeeds and supplies an effective positive lower bound on
$|y|$. Rational approximations to $1/y$, and hence $x/y$, follow with
computable error bounds. All four constructions use only rational arithmetic.
\end{proof}

\begin{proposition}[Real Algebraic Closure]
Let $p\in\Rc[X]$ be a nonzero polynomial. Every real root of $p$ is in
$\Rc$. Equivalently, the field $\Rc$ is real closed.
\end{proposition}
The real-root statement is the effective real-algebraic closure property;
it is recorded here without reproducing the considerably longer root-isolation
argument. It should not be confused with algebraic closure over $\C$:
$\Rc$ is not an algebraically closed field.

\begin{proposition}[Computability of $\pi$ and $e$]
The real numbers $\pi$ and $e$ are computable.
\end{proposition}
\begin{proof}
The alternating arctangent series gives
\[
 \frac{\pi}{4}=\arctan 1=\sum_{k=0}^{\infty}\frac{(-1)^k}{2k+1}.
\]
For $m\in\Nzero$ put
\[
 p_m=4\sum_{k=0}^{m}\frac{(-1)^k}{2k+1}\in\Q.
\]
The alternating-series estimate yields
\[
 |\pi-p_m|\le\frac{4}{2m+3}.
\]
Choosing, for instance, $m=2^{n+2}$ gives
$|\pi-p_m|<2^{-n}$ for every $n\in\Nzero$. The finite sum $p_m$ is
computable from $n$.

Similarly, let $e_m=\sum_{k=0}^{m}1/k!$. For $m\ge1$,
\[
 0<e-e_m=\sum_{k=m+1}^{\infty}\frac1{k!}
 \le\frac{2}{(m+1)!}.
\]
Taking $m=n+3$ makes the error strictly smaller than $2^{-n}$.
Thus $e$ is computable as well.
\end{proof}

\begin{proposition}[Countability of the Computable Reals]
The set $\Rc$ is countably infinite. Consequently, there exist
noncomputable real numbers.
\end{proposition}
\begin{proof}
There are countably many finite descriptions of algorithms, and every
computable real has at least one such description. Hence $\Rc$ is at most
countable. On the other hand, every rational number is computable via its
constant approximation sequence, so $\Q\subseteq\Rc$ and $\Rc$ is infinite.
Since $\R$ is uncountable, $\R\setminus\Rc\ne\varnothing$.
\end{proof}

\begin{fact}[Hilbert's Tenth Problem]
There is no algorithm that decides, for every multivariate polynomial with
integer coefficients, whether it has an integer zero. In particular, the set
of integer polynomials admitting an integer solution is not decidable.
\end{fact}

\begin{proposition}[A Noncomputable Real]
Let $\mathcal P$ be the set of all integer-coefficient polynomials in finitely
many variables, and let
$\phi\colon\mathcal P\to\Nzero$ be a computable bijective encoding.
Write $\mathcal P_{\mathrm{root}}$ for the polynomials having an integer zero.
Then
\[
 r:=\sum_{p\in\mathcal P}
 2^{-(\phi(p)+1)}\,\mathds{1}_{\mathcal P_{\mathrm{root}}}(p)
\]
is a noncomputable real in $[0,1]$.
\end{proposition}
\begin{proof}
The binary digit at position $\phi(p)$ is $1$ precisely when
$p\in\mathcal P_{\mathrm{root}}$. Both this set and its complement are infinite;
for example, among polynomials of one variable there are infinitely many
with an integer zero and infinitely many without one. Thus the binary
sequence has infinitely many zeros and ones, so the expansion is neither
eventually zero nor eventually one and $r$ is not dyadic.

If $r$ were computable, the equivalence theorem would give a computable
binary expansion. Because $r$ is not dyadic, its binary expansion is unique;
its digit at position $\phi(p)$ could be computed from $p$. This would decide
whether $p$ has an integer zero, contradicting Hilbert's Tenth Problem.
\end{proof}

\begin{remark}
More generally, any undecidable subset of $\Nzero$ can be encoded as the
binary digits of a noncomputable real, provided its characteristic sequence
is not eventually constant. Thus undecidability supplies explicit examples
that cardinality alone does not identify.
\end{remark}

\section{Computable Functions}

The definition of a computable function must work with approximations to an
unknown real input, rather than requiring an exact representation of that input.

\begin{definition}[Computable Function]
Let $D\subseteq\R$. A function $f\colon D\to\R$ is \emph{computable} if
there exists an oracle algorithm $\A$ with the following property. For every
$x\in D$, every function $\alpha\colon\Nzero\to\Q$ satisfying
\[
 |\alpha(i)-x|<2^{-i}\qquad(i\in\Nzero),
\]
and every precision $n\in\Nzero$, the algorithm $\A$, given access to
$\alpha$ and the input $n$, halts and returns a rational number
$\A^{\alpha}(n)$ such that
\[
 |\A^{\alpha}(n)-f(x)|<2^{-n}.
\]
The algorithm may query finitely many values $\alpha(i)$ during a computation.
The definition extends componentwise to functions between subsets of
$\R^k$ and $\R^m$.
\end{definition}

\begin{remark}
The input name $\alpha$ need not itself be computable: the algorithm must
work for \emph{every} sufficiently accurate rational name of $x$, including
names of noncomputable reals. What must be computable is the procedure
transforming such a name into approximations to $f(x)$.
\end{remark}

If $x\in\Rc$, choosing a computable name $\alpha$ and executing $\A$ gives
computable rational approximations to $f(x)$. Thus a computable function sends
computable inputs to computable outputs.

\begin{example}[Basic Computable Functions]
For each $r\in\Rc$, the maps $x\mapsto x+r$ and $x\mapsto rx$ are
computable. Other standard examples include
\[
 x\mapsto |x|,\quad x\mapsto e^x,\quad x\mapsto\sin x,\quad
 x\mapsto\sqrt{x},\quad x\mapsto\log x,
\]
with the natural domains $\R$, $\R$, $\R$, $[0,\infty)$, and
$(0,\infty)$, respectively. The exponential function is treated
constructively below.

For functions of several real variables, addition, subtraction, multiplication,
$\min$ and $\max$ are computable; so is division on the domain where the
denominator is nonzero. In each case, rational input approximations suffice
to control the output error.
\end{example}

\begin{theorem}[Closure under Composition]
Let $f\colon D\to E$ and $g\colon E\to\R^m$ be computable, where
$D\subseteq\R^k$ and $E\subseteq\R^\ell$. Then $g\circ f$ is computable.
\end{theorem}
\begin{proof}
Given a name of $x\in D$, the algorithm for $f$ produces, on request,
rational vectors $u_j$ with $\|u_j-f(x)\|_\infty<2^{-(j+2)}$.
The map $j\mapsto u_j$ is therefore a valid name of $f(x)$, obtained
on demand from the original input name. Feed this name to the algorithm for
$g$. Each computation of $g$ makes only finitely many requests, and each
request to $f$ terminates. The resulting rational output approximates
$g(f(x))$ to the requested precision.
\end{proof}

\begin{theorem}[Computability Implies Continuity]
Every computable function $f\colon D\to\R^m$, with
$D\subseteq\R^k$, is continuous on $D$ (with the relative topology).
\end{theorem}
\begin{proof}
Fix $x\in D$ and $\varepsilon>0$. Choose $n\in\Nzero$ with
$2^{1-n}<\varepsilon$. Select a rational name $\alpha$ of $x$ having
extra precision, namely
\[
 \|\alpha(i)-x\|_\infty<2^{-(i+2)}
 \qquad(i\in\Nzero).
\]
Run the algorithm for $f$ on $\alpha$ at precision $n$. As it halts, it
queries only finitely many indices, say those in a finite set $F$.
There exists $\delta>0$ such that whenever $y\in D$ satisfies
$\|x-y\|_\infty<\delta$, the same rational values $\alpha(i)$,
$i\in F$, are also valid approximations to $y$ with error less than
$2^{-i}$. Extend these finitely many values to a rational name $\beta$
of $y$. The algorithm receives exactly the same answers to all queries
and therefore makes the same computation, returning the same rational output
$q$. Consequently,
\[
 \|f(x)-f(y)\|_\infty
 \le \|f(x)-q\|_\infty+\|q-f(y)\|_\infty
 <2^{1-n}<\varepsilon.
\]
Thus $f$ is continuous at $x$.
\end{proof}

\begin{remark}
Continuity is necessary, not sufficient, for computability: if
$c\in\R\setminus\Rc$, the constant function $x\mapsto c$ is continuous
but not computable. In particular, a function with a discontinuity at a
point of its domain cannot be computable there.
\end{remark}

\begin{example}[Discontinuous Functions and Their Domains]
The step function $x\mapsto\mathds{1}_{\{x>0\}}$ is not computable on
$\R$, since it is discontinuous at $0$. Its restriction to
$\R\setminus\{0\}$ is computable: approximate the nonzero input until
its sign can be certified, and return the corresponding bit.
The same principle applies to the strict-comparison function on the domain
$x\ne y$, as shown below. Allowing set-valued outputs is another possible
way to treat boundary ambiguity, but requires a separate definition of
computability and is not pursued here.
\end{example}

\begin{theorem}[Computability of the Exponential Function]
The restriction $\exp\colon[-\tfrac12,\tfrac12]\to\R$ is computable.
\end{theorem}
\begin{proof}
Let $x\in[-\tfrac12,\tfrac12]$, let $\alpha$ be any rational name of
$x$, and fix $n\in\Nzero$. Put $m=n+4$, $d=n+4$, and
\[
 \widehat{x}=\alpha(m)\in\Q,\qquad
 P_d(t)=\sum_{j=0}^{d}\frac{t^j}{j!}\in\Q[t].
\]
The algorithm returns $P_d(\widehat{x})$, which is rational and can be
computed using finite rational arithmetic. Since
$|\widehat{x}-x|<2^{-m}\le1/16$, we have $|\widehat{x}|<1$.

By the mean value theorem and $e<3$, the error due to approximating the
input satisfies
\[
 |e^x-e^{\widehat{x}}|
 \le 3|x-\widehat{x}|<3\cdot 2^{-m}
 <2^{-(n+1)}.
\]
Meanwhile, the Taylor remainder obeys
\[
 |e^{\widehat{x}}-P_d(\widehat{x})|
 \le\sum_{j=d+1}^{\infty}\frac{|\widehat{x}|^j}{j!}
 <\sum_{j=d+1}^{\infty}\frac1{j!}
 \le\frac{2}{(d+1)!}<2^{-(n+1)}.
\]
The triangle inequality therefore gives
$|e^x-P_d(\widehat{x})|<2^{-n}$, as required.
\end{proof}

For arbitrary real $x$, one can effectively choose $k\in\Nzero$ such
that $|x/2^k|\le\tfrac12$, using a coarse rational approximation to $x$.
The identity $e^x=(e^{x/2^k})^{2^k}$ and closure under multiplication then
extend computability of the exponential function to all of $\R$.

\begin{theorem}[Noncomputability of Strict Comparison]
The strict-comparison function
\[
 L\colon\R^2\to\{0,1\},\qquad
 L(x,y)=\begin{cases}1,&x<y,\\0,&x\ge y,\end{cases}
\]
is not computable.
\end{theorem}
\begin{proof}
Suppose an algorithm computes $L$. Because its values lie in $\{0,1\}$,
requesting an approximation with error below $1/4$ and testing whether the
result exceeds $1/2$ yields an algorithm returning the exact bit.
Give this decision algorithm the two names
\[
 \alpha(i)=0,\qquad \beta(i)=2^{-(i+1)}\qquad(i\in\Nzero),
\]
both of which represent the real number $0$. Hence the algorithm must
halt and output $0$. During this run, only finitely many indices of
$\beta$ are queried; let $N\in\Nzero$ exceed or equal all of them.
Define another rational name by
\[
 \beta'(i)=\begin{cases}
    2^{-(i+1)},&i\le N,\\
    2^{-(N+1)},&i>N.
 \end{cases}
\]
For $i\le N$, its distance from $2^{-(N+1)}$ is strictly smaller than
$2^{-i}$, and for $i>N$ the distance is zero. Thus $\beta'$ is a valid
name of $y'=2^{-(N+1)}>0$.
The algorithm cannot distinguish $(\alpha,\beta')$ from
$(\alpha,\beta)$ during its finite computation; it still returns $0$.
But $L(0,y')=1$, a contradiction.
\end{proof}

\begin{proposition}[Comparison of Unequal Real Numbers]
The restriction
\[
 L\colon\{(x,y)\in\R^2:x\ne y\}\to\{0,1\}
\]
of strict comparison is computable.
\end{proposition}
\begin{proof}
Given names $\alpha,\beta$ of $x,y$, inspect successively the rational
intervals
\[
 I_n=(\alpha(n)-2^{-n},\,\alpha(n)+2^{-n}),\qquad
 J_n=(\beta(n)-2^{-n},\,\beta(n)+2^{-n}).
\]
If $\sup I_n<\inf J_n$, output $1$; if
$\sup J_n<\inf I_n$, output $0$. These comparisons are rational and
hence decidable. Since $x\ne y$ and the intervals shrink to their
respective points, eventually one of the tests succeeds. The sign of
$x-y$ is then certified, so the answer is correct.
\end{proof}

\begin{remark}
The obstruction to comparison is equality, not numerical closeness.
Without a promise that $x\ne y$, no finite amount of approximation data
can always determine whether two real inputs are equal. With the promise,
refining the approximations eventually separates them. The same distinction
applies to $\le$, $\ge$ and equality tests.
\end{remark}

\end{document}
